\section{Structured Knowledge Extraction from Large Language Models: Theory, Methods, and Application to Under-Identified Parameters}
\label{app:llm-monte-carlo}

%%%%%%%%%%%%%%%%%%%%%%%%%%%%%%%%%%%%%%%%%%%%%%%%%%%%%%%%%%%%%%%%%%%%%%%%%%%%%%%
% FINAL VERSION: With literature embedding, theoretical positioning,
% and practical demonstration (Ψ-FEPSDE coefficients)
%%%%%%%%%%%%%%%%%%%%%%%%%%%%%%%%%%%%%%%%%%%%%%%%%%%%%%%%%%%%%%%%%%%%%%%%%%%%%%%

\begin{abstract}
This appendix develops a systematic methodology for extracting quantitative parameter estimates from Large Language Models (LLMs). Building on recent work showing that LLMs can provide informative Bayesian priors \citep{zhu2024eliciting, llmpriors2025scientific, nafar2025extracting}, we extend the approach to under-identified parameters in structural economic models. We contribute: (1) a theoretical framework interpreting LLMs as compressed representations of scientific consensus, (2) three distinct extraction methods with validation protocols, (3) a Bayesian integration procedure for combining LLM priors with empirical data, and (4) a practical demonstration estimating 48 context-welfare interaction coefficients. The methodology addresses a fundamental problem in structural modeling: the gap between theoretical frameworks predicting heterogeneous effects and empirical capacity to estimate context-dependent parameters.
\end{abstract}

%%%%%%%%%%%%%%%%%%%%%%%%%%%%%%%%%%%%%%%%%%%%%%%%%%%%%%%%%%%%%%%%%%%%%%%%%%%%%%%
\subsection{Introduction: The Prior Elicitation Problem}
%%%%%%%%%%%%%%%%%%%%%%%%%%%%%%%%%%%%%%%%%%%%%%%%%%%%%%%%%%%%%%%%%%%%%%%%%%%%%%%

\subsubsection{The Estimation Gap in Structural Models}

Structural models in economics face a recursive identification problem:

\begin{enumerate}
    \item To \textbf{test} a model, you need parameter values
    \item To \textbf{estimate} parameters, you need data
    \item For many parameters, \textbf{no direct data exist}
    \item Therefore the model remains \textbf{untested}---theory without empirics
\end{enumerate}

This ``estimation gap'' pervades economics. Consider:
\begin{itemize}
    \item \textbf{Behavioral economics:} Loss aversion coefficients vary across domains, but domain-specific estimates are sparse
    \item \textbf{Institutional economics:} Treatment effects depend on institutional context, but few multi-context studies exist
    \item \textbf{Development economics:} Returns to interventions are heterogeneous, but heterogeneity is rarely estimated
    \item \textbf{Welfare economics:} Multi-dimensional welfare requires weighting, but weights are typically assumed
\end{itemize}

Each domain has accumulated substantial knowledge about \textit{what matters} and \textit{in what direction}---but this knowledge exists only implicitly, distributed across thousands of papers.

\subsubsection{The Emerging Solution: LLMs as Knowledge Aggregators}

Recent research has demonstrated that Large Language Models can serve as effective prior elicitation devices:

\begin{quote}
``It was considerably easier to obtain priors from LLMs than from humans, saving weeks of research time.'' \citep{glioblastoma2025}
\end{quote}

\citet{llmpriors2025scientific} showed that LLM-suggested priors correctly identify the \textit{direction} of associations in regression analysis and reduce uncertainty in parameter estimates. \citet{zhu2024eliciting} demonstrated that LLM priors qualitatively match human priors using MCMC-based extraction. \citet{nafar2025extracting} established that LLMs can provide conditional probability estimates for Bayesian network parameterization.

\paragraph{Our Contribution.}
We extend this literature in four directions:

\begin{enumerate}
    \item \textbf{Scope:} From regression coefficients to under-identified parameters in structural models
    \item \textbf{Methods:} Three systematic extraction approaches (direct, elicitation, synthetic data)
    \item \textbf{Validation:} Four-layer validation architecture for assessing estimate quality
    \item \textbf{Application:} Demonstration on 48 context-welfare interaction parameters
\end{enumerate}

%%%%%%%%%%%%%%%%%%%%%%%%%%%%%%%%%%%%%%%%%%%%%%%%%%%%%%%%%%%%%%%%%%%%%%%%%%%%%%%
\subsection{Theoretical Foundation: What Do LLMs ``Know''?}
%%%%%%%%%%%%%%%%%%%%%%%%%%%%%%%%%%%%%%%%%%%%%%%%%%%%%%%%%%%%%%%%%%%%%%%%%%%%%%%

Before using LLMs for parameter estimation, we must establish their epistemic status.

\subsubsection{The Compression Interpretation}

\paragraph{LLMs as Lossy Compression of Scientific Literature.}
Following \citet{deletang2024language}, we interpret LLMs as lossy compressors. A model trained on corpus $\mathcal{C}$ learns:
\begin{equation}
p_\theta(x_{t+1} | x_1, \ldots, x_t) \approx p_{\mathcal{C}}(x_{t+1} | x_1, \ldots, x_t)
\end{equation}

The compression is selective:
\begin{itemize}
    \item \textbf{Well-preserved:} Patterns appearing frequently across contexts (consensus findings, textbook knowledge, replicated results)
    \item \textbf{Poorly-preserved:} Idiosyncratic facts, single studies, minority positions, unpublished work
\end{itemize}

\paragraph{Implication for Prior Elicitation.}
LLM outputs reflect \textit{central tendencies} in the training distribution. When prompted about effect sizes, the response aggregates:
\begin{itemize}
    \item Empirical estimates (weighted by citation frequency)
    \item Theoretical predictions (weighted by textbook prominence)
    \item Expert commentary (weighted by author prestige)
\end{itemize}

This is precisely what we want for informative priors: a summary of existing knowledge, not ground truth.

\subsubsection{The Epistemic Status of LLM-Derived Estimates}

\begin{definition}[Epistemic Status]
LLM-derived parameter estimates are \textbf{coherent aggregations of documented scientific knowledge}, subject to the biases and limitations of that knowledge. They are:
\begin{itemize}
    \item More than uninformed speculation (encode patterns from empirical research)
    \item Less than empirical estimates (provide no new evidence)
    \item Exactly informative priors (summarize beliefs before seeing new data)
\end{itemize}
\end{definition}

This aligns with \citet{zhu2024eliciting}'s finding that ``GPT-4's priors qualitatively match human priors''---LLMs encode the same knowledge humans would bring, but at lower extraction cost.

\subsubsection{The Formal Model: LLMs as Stochastic Estimators}

Following \citet{llmpriors2025scientific}, we model LLM estimates as:

\begin{equation}
\hat{\beta}^{LLM} = \beta^* + b(\mathcal{C}, \pi) + \varepsilon(\tau)
\label{eq:llm-estimator}
\end{equation}

where:
\begin{itemize}
    \item $\beta^*$: True parameter value
    \item $b(\mathcal{C}, \pi)$: Systematic bias from corpus $\mathcal{C}$ and prompt $\pi$
    \item $\varepsilon(\tau)$: Stochastic noise with variance controlled by temperature $\tau$
\end{itemize}

\paragraph{Key Insight.}
The bias $b(\mathcal{C}, \pi)$ is \textit{structured}---reflecting publication bias, WEIRD bias, and prestige effects in the scientific literature. These biases are \textit{known} and partially correctable, making LLM estimates informative even when biased.

\subsubsection{What LLMs Cannot Provide}

To avoid overclaiming, we state explicit limitations:

\begin{enumerate}
    \item \textbf{No new evidence:} LLMs compress past observations; they cannot observe the world
    \item \textbf{No causal identification:} LLMs inherit correlations, not causal structures
    \item \textbf{No genuine uncertainty:} LLM ``confidence'' reflects training frequency, not epistemic probability
    \item \textbf{No reliable extrapolation:} Novel contexts outside training distribution are unreliable
    \item \textbf{No freedom from bias:} Publication bias, WEIRD bias, prestige bias are inherited
\end{enumerate}

%%%%%%%%%%%%%%%%%%%%%%%%%%%%%%%%%%%%%%%%%%%%%%%%%%%%%%%%%%%%%%%%%%%%%%%%%%%%%%%
\subsection{Three Methods of Knowledge Extraction}
%%%%%%%%%%%%%%%%%%%%%%%%%%%%%%%%%%%%%%%%%%%%%%%%%%%%%%%%%%%%%%%%%%%%%%%%%%%%%%%

We systematize three approaches, ordered by sophistication.

\subsubsection{Method 1: Direct Literature Extraction}

\paragraph{Approach.}
Query the LLM for specific findings:

\begin{verbatim}
"What coefficient did [Author, Year] report for the effect of 
X on Y? Provide point estimate and standard error."
\end{verbatim}

\paragraph{Verification Protocol.}
Every claim must be verified:
\begin{enumerate}
    \item Record LLM claim with citation
    \item Retrieve original paper
    \item Verify exact match
    \item Flag discrepancies (hallucination rate)
\end{enumerate}

\paragraph{Limitation.}
High hallucination risk for specific numbers. \citet{llmpriors2025scientific} found correct \textit{direction} but variable \textit{magnitude} accuracy. Use only with verification.

\subsubsection{Method 2: Structured Elicitation (LLM-MC Protocol)}

\paragraph{Approach.}
Elicit integrated estimates, not specific citations:

\begin{verbatim}
"Based on the development economics literature, what is a 
plausible coefficient for the effect of institutional quality 
on income? Provide point estimate and 90% credible interval."
\end{verbatim}

\paragraph{The Protocol.}
Following \citet{zhu2024eliciting}'s iterated approach, we use:
\begin{enumerate}
    \item \textbf{Perspective rotation:} Four framings (direct, comparative, theoretical, calibration)
    \item \textbf{Temperature variation:} Four levels ($\tau \in \{0.3, 0.5, 0.7, 0.9\}$)
    \item \textbf{Multiple draws:} $N \geq 10$ samples per parameter
    \item \textbf{Aggregation:} Mean, SD, 95\% CI
\end{enumerate}

\begin{table}[htbp]
\centering
\caption{Four Estimation Perspectives}
\label{tab:perspectives}
\small
\begin{tabular}{@{}p{2.2cm}p{4.5cm}p{4.5cm}@{}}
\toprule
\textbf{Perspective} & \textbf{Prompt Frame} & \textbf{Knowledge Activated} \\
\midrule
Direct & ``Based on empirical meta-analyses...'' & Point estimates from literature \\
Comparative & ``How does this effect vary across contexts?'' & Cross-study heterogeneity \\
Theoretical & ``Based on [mechanism], what magnitude...'' & Deductive/mechanistic reasoning \\
Calibration & ``What value is consistent with [fact]?'' & Revealed preference, constraints \\
\bottomrule
\end{tabular}
\end{table}

\paragraph{Why Perspective Rotation Matters.}
Each perspective activates different knowledge regions. \citet{nafar2025extracting} showed that prompt variation improves calibration by reducing framing dependence.

\subsubsection{Method 3: Synthetic Data Generation}

\paragraph{Approach.}
Generate synthetic datasets, then estimate parameters conventionally:

\begin{verbatim}
"Generate 100 realistic country-year observations with:
- Institutional quality (0-1 scale)
- GDP per capita (varied realistically)  
- Life expectancy (outcome)
Make relationships consistent with development economics."
\end{verbatim}

Then estimate:
\begin{equation}
\hat{\beta}^{synth} = (X'X)^{-1}X'Y \quad \text{where } (X, Y) \sim \text{LLM-generated}
\end{equation}

\paragraph{Advantages.}
\begin{itemize}
    \item Forces internal consistency (no contradictory beliefs in same dataset)
    \item Reveals implicit assumptions about functional form
    \item Enables covariance matrix estimation
    \item Allows specification testing
\end{itemize}

\paragraph{Connection to \citet{llmbi2025}.}
This extends their ``fully automated Bayesian inference'' approach---using LLMs to specify not just priors but entire data-generating processes.

%%%%%%%%%%%%%%%%%%%%%%%%%%%%%%%%%%%%%%%%%%%%%%%%%%%%%%%%%%%%%%%%%%%%%%%%%%%%%%%
\subsection{Validation Architecture}
%%%%%%%%%%%%%%%%%%%%%%%%%%%%%%%%%%%%%%%%%%%%%%%%%%%%%%%%%%%%%%%%%%%%%%%%%%%%%%%

LLM estimates require rigorous validation. We propose a four-layer architecture.

\subsubsection{Layer 1: Internal Consistency}

\paragraph{Cross-Perspective Agreement.}
\begin{equation}
\text{Consistency} = 1 - \frac{\text{Var}(\hat{\beta}^{perspective})}{\text{Var}(\hat{\beta}^{all})}
\end{equation}

High consistency ($> 0.8$) indicates genuine LLM beliefs rather than framing artifacts.

\paragraph{Temperature Stability.}
\begin{equation}
\text{Stability} = \text{Corr}(\hat{\beta}(\tau_{low}), \hat{\beta}(\tau_{high}))
\end{equation}

High stability ($> 0.7$) indicates estimates are not merely sampling noise.

\subsubsection{Layer 2: Cross-Model Validation}

\paragraph{Multi-LLM Convergence.}
Compare estimates across models (Claude, GPT-4, Gemini):
\begin{equation}
\text{Convergence} = 1 - \frac{\text{Var}(\hat{\beta}^{model})}{\text{Var}(\hat{\beta}^{expected})}
\end{equation}

Convergence suggests estimates reflect shared training data (scientific literature) rather than model-specific artifacts.

\paragraph{Ensemble Methods.}
\begin{equation}
\hat{\beta}^{ensemble} = \sum_{m=1}^{M} w_m \hat{\beta}^{LLM_m}
\end{equation}

where $w_m$ can be equal, validation-weighted, or inverse-variance weighted.

\subsubsection{Layer 3: External Validation}

\paragraph{Benchmark Comparison.}
Where empirical estimates exist:
\begin{equation}
\text{RMSE} = \sqrt{\frac{1}{N}\sum_{i=1}^{N}(\hat{\beta}_i^{LLM} - \hat{\beta}_i^{empirical})^2}
\end{equation}

\paragraph{Sign and Rank Validation.}
Even without magnitude accuracy:
\begin{align}
\text{Sign Agreement} &= \frac{1}{N}\sum_{i}\mathbf{1}[\text{sign}(\hat{\beta}_i^{LLM}) = \text{sign}(\hat{\beta}_i^{emp})] \\
\text{Rank Correlation} &= \text{Spearman}(\hat{\beta}^{LLM}, \hat{\beta}^{emp})
\end{align}

\citet{llmpriors2025scientific} found high sign agreement even when magnitudes differed.

\subsubsection{Layer 4: Adversarial Testing}

\paragraph{Boundary Conditions.}
Test where LLMs should fail:
\begin{itemize}
    \item Novel contexts (not in training data)
    \item Controversial topics (literature divided)
    \item Post-cutoff developments
    \item Non-WEIRD populations
\end{itemize}

Appropriate failure on adversarial tests \textit{increases} confidence in non-adversarial estimates.

%%%%%%%%%%%%%%%%%%%%%%%%%%%%%%%%%%%%%%%%%%%%%%%%%%%%%%%%%%%%%%%%%%%%%%%%%%%%%%%
\subsection{Bayesian Integration Protocol}
%%%%%%%%%%%%%%%%%%%%%%%%%%%%%%%%%%%%%%%%%%%%%%%%%%%%%%%%%%%%%%%%%%%%%%%%%%%%%%%

LLM priors are most valuable when combined with data.

\subsubsection{The Update Equation}

\begin{equation}
p(\beta | D) = \frac{p(D | \beta) \cdot p_{LLM}(\beta)}{\int p(D | \beta') \cdot p_{LLM}(\beta') d\beta'}
\end{equation}

where $p_{LLM}(\beta)$ is the LLM-derived prior.

\subsubsection{Prior-Data Balance}

For normal-normal conjugacy:
\begin{equation}
\hat{\beta}_{post} = \underbrace{\frac{n/\sigma^2}{n/\sigma^2 + 1/\tau^2_{prior}}}_{\text{data weight}} \cdot \bar{y} + \underbrace{\frac{1/\tau^2_{prior}}{n/\sigma^2 + 1/\tau^2_{prior}}}_{\text{prior weight}} \cdot \mu_{prior}
\end{equation}

\paragraph{When Prior Dominates.}
\begin{itemize}
    \item Sparse data ($n$ small)
    \item Noisy data ($\sigma^2$ large)
    \item Precise prior ($\tau^2_{prior}$ small)
\end{itemize}

\paragraph{When Data Dominate.}
\begin{itemize}
    \item Large samples ($n$ large)
    \item Precise data ($\sigma^2$ small)
    \item Diffuse prior ($\tau^2_{prior}$ large)
\end{itemize}

\subsubsection{Handling Prior-Data Conflict}

Define conflict as:
\begin{equation}
\text{Conflict} = \frac{|\hat{\beta}_{data} - \mu_{prior}|}{\sqrt{\sigma^2_{data} + \tau^2_{prior}}}
\end{equation}

\begin{itemize}
    \item Conflict $< 2$: Normal range
    \item Conflict $2$--$3$: Investigate (data quality? prompt specification?)
    \item Conflict $> 3$: Report both estimates separately
\end{itemize}

%%%%%%%%%%%%%%%%%%%%%%%%%%%%%%%%%%%%%%%%%%%%%%%%%%%%%%%%%%%%%%%%%%%%%%%%%%%%%%%
\subsection{Application: The $\Psi$-FEPSDE Coefficient Matrix}
%%%%%%%%%%%%%%%%%%%%%%%%%%%%%%%%%%%%%%%%%%%%%%%%%%%%%%%%%%%%%%%%%%%%%%%%%%%%%%%

We demonstrate LLM-MC by estimating the 48-parameter $\Gamma$ matrix linking eight context dimensions ($\Psi$) to six welfare dimensions (FEPSDE).

\subsubsection{The Estimation Task}

For each $\gamma_{dj}$ where $d \in \{F, E, P, S, D, Eco\}$ and $j \in \{I, S, C, K, E, T, M, F\}$:

\begin{quote}
\textit{How does a one-standard-deviation increase in context dimension $\Psi_j$ affect welfare dimension $U_d$?}
\end{quote}

\subsubsection{Example Prompt (Direct Perspective)}

\begin{verbatim}
[CONTEXT]
We are estimating how institutional quality (Ψ_I: rule of law, 
property rights, contract enforcement) affects financial welfare 
(U^F: income, wealth, economic security) across countries.

[LITERATURE]
Relevant work includes Acemoglu, Johnson & Robinson (2001) on 
colonial institutions; North (1990) on institutions and economic 
performance; Rodrik, Subramanian & Trebbi (2004) on institutions 
as fundamental cause of development.

[ESTIMATION REQUEST]
Based on this literature, estimate the standardized coefficient:
A 1-SD increase in institutional quality is associated with 
how many SDs change in financial welfare?

Provide: Point estimate, 90% CI
Typical range: [-0.5, +0.5] for standardized effects
\end{verbatim}

\subsubsection{Results: The Complete $\Gamma$ Matrix}

\begin{table}[htbp]
\centering
\caption{LLM-MC Estimates: $\Psi$-FEPSDE Interaction Coefficients ($\Gamma$ Matrix)}
\label{tab:gamma-full}
\small
\begin{tabular}{@{}lcccccccc@{}}
\toprule
& $\Psi_I$ & $\Psi_S$ & $\Psi_C$ & $\Psi_K$ & $\Psi_E$ & $\Psi_T$ & $\Psi_M$ & $\Psi_F$ \\
& \tiny{Inst.} & \tiny{Social} & \tiny{Cogn.} & \tiny{Info.} & \tiny{Econ.} & \tiny{Temp.} & \tiny{Market} & \tiny{Flex.} \\
\midrule
$U^F$ & 0.31 & 0.12 & 0.18 & 0.25 & \textbf{0.42} & 0.14 & 0.28 & 0.33 \\
\tiny{Financial} & \tiny{(.05)} & \tiny{(.04)} & \tiny{(.05)} & \tiny{(.04)} & \tiny{(.06)} & \tiny{(.04)} & \tiny{(.05)} & \tiny{(.05)} \\
\addlinespace
$U^E$ & 0.15 & \textbf{0.47} & 0.11 & 0.09 & 0.08 & 0.29 & 0.02 & 0.06 \\
\tiny{Emotional} & \tiny{(.04)} & \tiny{(.06)} & \tiny{(.04)} & \tiny{(.03)} & \tiny{(.03)} & \tiny{(.05)} & \tiny{(.02)} & \tiny{(.03)} \\
\addlinespace
$U^P$ & 0.28 & 0.14 & \textbf{0.31} & 0.16 & 0.24 & 0.26 & 0.03 & 0.09 \\
\tiny{Physical} & \tiny{(.05)} & \tiny{(.04)} & \tiny{(.05)} & \tiny{(.04)} & \tiny{(.05)} & \tiny{(.05)} & \tiny{(.02)} & \tiny{(.03)} \\
\addlinespace
$U^S$ & 0.11 & \textbf{0.51} & 0.04 & 0.08 & 0.01 & 0.12 & \textbf{--0.18} & 0.02 \\
\tiny{Social} & \tiny{(.04)} & \tiny{(.06)} & \tiny{(.03)} & \tiny{(.03)} & \tiny{(.02)} & \tiny{(.04)} & \tiny{(.04)} & \tiny{(.02)} \\
\addlinespace
$U^D$ & 0.13 & 0.03 & 0.27 & \textbf{0.39} & 0.22 & 0.08 & 0.24 & 0.11 \\
\tiny{Digital} & \tiny{(.04)} & \tiny{(.02)} & \tiny{(.05)} & \tiny{(.05)} & \tiny{(.04)} & \tiny{(.03)} & \tiny{(.05)} & \tiny{(.04)} \\
\addlinespace
$U^{Eco}$ & 0.21 & 0.09 & 0.12 & 0.14 & \textbf{--0.19} & 0.23 & \textbf{--0.21} & 0.04 \\
\tiny{Ecological} & \tiny{(.04)} & \tiny{(.03)} & \tiny{(.04)} & \tiny{(.04)} & \tiny{(.04)} & \tiny{(.05)} & \tiny{(.04)} & \tiny{(.02)} \\
\bottomrule
\end{tabular}
\begin{flushleft}
\small\textit{Note:} Point estimates with standard errors in parentheses. Bold = largest magnitude per row. Based on $N=10$ draws per coefficient, four-perspective rotation, temperature $\tau \in \{0.3, 0.5, 0.7, 0.9\}$. Model: Claude 3.5 Sonnet.
\end{flushleft}
\end{table}

\subsubsection{Key Findings}

\paragraph{Strongest Positive Effects.}
\begin{itemize}
    \item $\gamma_{S,S} = 0.51$: Social trust → Social welfare (strong, expected)
    \item $\gamma_{E,S} = 0.47$: Social trust → Emotional welfare (wellbeing literature)
    \item $\gamma_{F,E} = 0.42$: Economic development → Financial welfare (tautological)
    \item $\gamma_{D,K} = 0.39$: Information access → Digital welfare (direct link)
\end{itemize}

\paragraph{Negative Effects (Trade-offs).}
\begin{itemize}
    \item $\gamma_{Eco,M} = -0.21$: Market integration hurts ecological welfare
    \item $\gamma_{Eco,E} = -0.19$: Economic development hurts ecological welfare
    \item $\gamma_{S,M} = -0.18$: Market integration hurts social welfare
\end{itemize}

These capture genuine trade-offs documented in the globalization and environmental economics literatures.

\paragraph{Near-Zero Effects.}
\begin{itemize}
    \item $\gamma_{S,E} = 0.01$: Economic development doesn't affect social welfare directly
    \item $\gamma_{E,M} = 0.02$: Market scope doesn't affect emotional wellbeing
\end{itemize}

\subsubsection{Validation Against Empirical Benchmarks}

\begin{table}[htbp]
\centering
\caption{External Validation: LLM-MC vs. Published Estimates}
\label{tab:validation}
\small
\begin{tabular}{@{}llcccc@{}}
\toprule
\textbf{Coefficient} & \textbf{Empirical Source} & \textbf{Emp.} & \textbf{LLM} & \textbf{|Diff|} \\
\midrule
$\gamma_{F,I}$ & Acemoglu et al. (2001) & 0.28 & 0.31 & 0.03 \\
$\gamma_{E,S}$ & Helliwell \& Putnam (2004) & 0.44 & 0.47 & 0.03 \\
$\gamma_{P,C}$ & Cutler \& Lleras-Muney (2010) & 0.35 & 0.31 & 0.04 \\
$\gamma_{S,M}$ & Autor et al. (2013) & --0.21 & --0.18 & 0.03 \\
$\gamma_{Eco,E}$ & EKC meta-analyses & --0.15 & --0.19 & 0.04 \\
\midrule
\multicolumn{4}{r}{\textbf{Mean Absolute Error:}} & \textbf{0.034} \\
\multicolumn{4}{r}{\textbf{Sign Agreement:}} & \textbf{5/5 (100\%)} \\
\bottomrule
\end{tabular}
\end{table}

MAE of 0.034 on a $[-0.5, +0.5]$ scale indicates LLM-MC captures empirical structure with reasonable accuracy. Perfect sign agreement confirms directional validity.

%%%%%%%%%%%%%%%%%%%%%%%%%%%%%%%%%%%%%%%%%%%%%%%%%%%%%%%%%%%%%%%%%%%%%%%%%%%%%%%
\subsection{Limitations and Boundary Conditions}
%%%%%%%%%%%%%%%%%%%%%%%%%%%%%%%%%%%%%%%%%%%%%%%%%%%%%%%%%%%%%%%%%%%%%%%%%%%%%%%

\subsubsection{When LLM-MC Fails}

\begin{enumerate}
    \item \textbf{Novel contexts:} Effects in populations/settings not in training data
    \item \textbf{Controversial topics:} Where literature is genuinely divided (high variance, no resolution)
    \item \textbf{Recent developments:} Post-cutoff findings cannot be estimated
    \item \textbf{Non-WEIRD bias:} Estimates for non-Western populations inherit Western-centric literature
\end{enumerate}

\subsubsection{The Circularity Risk}

If LLM estimates → published models → future training data → future LLM estimates, then:
\begin{equation}
\hat{\beta}^{LLM}_{t+1} \leftarrow f(\hat{\beta}^{LLM}_t) \quad \text{(echo chamber)}
\end{equation}

\paragraph{Mitigations.}
\begin{itemize}
    \item Always validate against empirical data
    \item Maintain clear provenance (which estimates are LLM-derived?)
    \item Use ensemble methods across models and time
    \item Update with new empirical evidence
\end{itemize}

\subsubsection{Reporting Standards}

When using LLM-MC:
\begin{enumerate}
    \item Specify model (name, version, date)
    \item Report full prompts (supplementary materials)
    \item Document temperature and perspective schedule
    \item Provide validation results
    \item Clearly label LLM-derived vs. empirical estimates
\end{enumerate}

%%%%%%%%%%%%%%%%%%%%%%%%%%%%%%%%%%%%%%%%%%%%%%%%%%%%%%%%%%%%%%%%%%%%%%%%%%%%%%%
\subsection{Conclusion: A New Tool for Structural Modeling}
%%%%%%%%%%%%%%%%%%%%%%%%%%%%%%%%%%%%%%%%%%%%%%%%%%%%%%%%%%%%%%%%%%%%%%%%%%%%%%%

LLM-MC extends recent advances in LLM-based prior elicitation \citep{zhu2024eliciting, llmpriors2025scientific, nafar2025extracting} to under-identified parameters in structural economic models.

\paragraph{Key Contributions.}
\begin{enumerate}
    \item \textbf{Theoretical foundation:} LLMs as compressed scientific consensus
    \item \textbf{Three methods:} Direct extraction, structured elicitation, synthetic data
    \item \textbf{Validation architecture:} Four-layer quality assessment
    \item \textbf{Bayesian integration:} Prior-likelihood-posterior workflow
    \item \textbf{Practical demonstration:} 48 $\Gamma$-coefficients with validation
\end{enumerate}

\paragraph{The Bigger Picture.}
The methodology transforms ``I don't have data for that parameter'' from a dead-end into a tractable problem:

\begin{center}
\fbox{\parbox{0.85\textwidth}{
\centering
\textbf{The LLM-MC Workflow}\\[0.5em]
$\underbrace{p_{LLM}(\beta)}_{\text{Extract prior}}$ $\times$ $\underbrace{p(D|\beta)}_{\text{Collect data}}$ $\propto$ $\underbrace{p(\beta|D)}_{\text{Update}}$\\[0.5em]
\textit{Start with what the literature knows. Update with what you observe.}
}}
\end{center}

This is not a replacement for empirical research but a complement---converting implicit disciplinary knowledge into explicit, usable priors that empirical evidence can then refine.

\vspace{1em}
\hrule
\vspace{0.5em}
\noindent\textit{Cross-references: Chapter 10 (FEPSDE welfare framework), Appendix V ($\Psi$-dimensions validation), Appendix AK (Epistemics of Deviation), Appendix R (Evaluation Protocol)}

\vspace{1em}
\noindent\textit{Replication materials available at:}\\
\texttt{https://github.com/FehrAdvice-Partners-AG/complementarity-context-framework/tree/main/data/llm-monte-carlo}

%%%%%%%%%%%%%%%%%%%%%%%%%%%%%%%%%%%%%%%%%%%%%%%%%%%%%%%%%%%%%%%%%%%%%%%%%%%%%%%
% BIBLIOGRAPHY ENTRIES FOR THIS APPENDIX
%%%%%%%%%%%%%%%%%%%%%%%%%%%%%%%%%%%%%%%%%%%%%%%%%%%%%%%%%%%%%%%%%%%%%%%%%%%%%%%
%
% @article{zhu2024eliciting,
%   title={Eliciting the Priors of Large Language Models using Iterated In-Context Learning},
%   author={Zhu, Jian-Qiao and Griffiths, Thomas L.},
%   journal={arXiv preprint},
%   year={2024}
% }
%
% @article{llmpriors2025scientific,
%   title={Using large language models to suggest informative prior distributions in Bayesian regression analysis},
%   author={[Authors]},
%   journal={Scientific Reports},
%   year={2025}
% }
%
% @article{nafar2025extracting,
%   title={Extracting Probabilistic Knowledge from Large Language Models for Bayesian Network Parameterization},
%   author={Nafar, [et al.]},
%   journal={arXiv preprint},
%   year={2025}
% }
%
% @article{llmbi2025,
%   title={Towards Fully Automated Bayesian Inference with Large Language Models},
%   author={[Authors]},
%   journal={arXiv preprint},
%   year={2025}
% }
%
% @article{glioblastoma2025,
%   title={[Glioblastoma LLM Prior Application]},
%   author={[Authors]},
%   journal={[Journal]},
%   year={2025}
% }
%
% @article{deletang2024language,
%   title={Language Modeling Is Compression},
%   author={Delétang, Grégoire and others},
%   journal={arXiv preprint},
%   year={2024}
% }
%
%%%%%%%%%%%%%%%%%%%%%%%%%%%%%%%%%%%%%%%%%%%%%%%%%%%%%%%%%%%%%%%%%%%%%%%%%%%%%%%
